\documentclass[11pt]{article}
\usepackage{amsmath,amssymb,amsthm}
\usepackage{algorithm}
\usepackage{algpseudocode}
\usepackage{bm}
\usepackage{hyperref}
\usepackage{enumitem}
\usepackage{geometry}
\geometry{margin=1in}
\newtheorem{theorem}{Theorem}
\newtheorem{lemma}{Lemma}
\newcommand{\E}{\mathbb{E}}
\newcommand{\R}{\mathbb{R}}
\newcommand{\norm}[1]{\left\lVert#1\right\rVert}
\newcommand{\grad}{\nabla}
\begin{document}

\section*{Differentially Private Stochastic Gradient Descent (DP‑SGD) with Gaussian Noise}

\section*{Mathematical Formulation}
Let $f:\R^{d}\rightarrow\R$ be the (possibly non‑convex) loss function we wish to minimise.
At iteration $t$ we maintain a parameter vector $x_t\in\R^{d}$.  
The algorithm proceeds as follows:

\begin{align}
\text{(1)  Gradient estimate (clipped)}\qquad &
\tilde g_t \;=\; \frac{1}{B}\sum_{i\in\mathcal{B}_t}
\operatorname{clip}\!\bigl(\grad f_i(x_t),C\bigr),                       \label{eq:clip}\\[4pt]
\text{(2)  Add Gaussian noise} \qquad &
\xi_t \;\sim\; \mathcal{N}\!\bigl(0,\sigma^{2}I_d\bigr),                     \label{eq:noise}\\[4pt]
\text{(3)  Noisy gradient} \qquad &
g_t \;=\; \tilde g_t + \xi_t,                                            \label{eq:noisygrad}\\[4pt]
\text{(4)  Parameter update} \qquad &
x_{t+1}\;=\;x_t - \eta_t g_t,                                            \label{eq:update}
\end{align}
where
\begin{itemize}
    \item $\mathcal{B}_t$ is a minibatch of size $B$ drawn uniformly at random,
    \item $\operatorname{clip}(v,C)=v\cdot\min\!\bigl\{1,\;C/\|v\|\bigr\}$ enforces $\|\operatorname{clip}(v,C)\|\le C$,
    \item $\eta_t>0$ is the learning rate,
    \item $\sigma>0$ controls the magnitude of the added Gaussian noise.
\end{itemize}
The privacy budget is tracked via the moments accountant (or Rényi DP); the only modification to standard SGD is the addition of the noise term $\xi_t$.

\section*{Pseudocode}
\begin{algorithm}[H]
\caption{DP‑SGD with Gaussian Noise}
\begin{algorithmic}[1]
\Require Initial point $x_0\in\R^{d}$, learning‑rate schedule $\{\eta_t\}_{t=0}^{T-1}$,
        clipping norm $C$, minibatch size $B$, noise scale $\sigma$,
        total iterations $T$.
\For{$t=0$ \textbf{to} $T-1$}
    \State Sample a minibatch $\mathcal{B}_t$ of size $B$.
    \State Compute per‑example gradients $\{ \grad f_i(x_t)\}_{i\in\mathcal{B}_t}$.
    \State Clip: $\tilde g_t \gets \frac{1}{B}\sum_{i\in\mathcal{B}_t}
            \operatorname{clip}\!\bigl(\grad f_i(x_t),C\bigr)$.
    \State Sample noise $\xi_t\sim\mathcal N(0,\sigma^{2}I_d)$.
    \State Form noisy gradient $g_t \gets \tilde g_t + \xi_t$.
    \State Update $x_{t+1} \gets x_t - \eta_t g_t$.
\EndFor
\State \Return $x_T$ (or an average $\bar x_T$).
\end{algorithmic}
\end{algorithm}

\section*{Dry Run Example}
Consider the simple scalar problem $f(w)=(w-3)^2$.
\begin{itemize}
    \item Gradient: $\grad f(w)=2(w-3)$.
    \item Initialise $w_0=0$, learning rate $\eta=0.1$, clipping $C=\infty$ (no clipping),
          and set the noise to zero for illustration ($\xi_t=0$).
\end{itemize}
We perform three iterations.

\begin{center}
\begin{tabular}{c|c|c|c}
Iteration $t$ & $\grad f(w_t)$ & $w_{t+1}=w_t-\eta\grad f(w_t)$ & $f(w_{t+1})$\\\hline
0 & $2(0-3)=-6$ & $0-0.1(-6)=0.6$ & $(0.6-3)^2=5.76$\\
1 & $2(0.6-3)=-4.8$ & $0.6-0.1(-4.8)=1.08$ & $(1.08-3)^2=3.6864$\\
2 & $2(1.08-3)=-3.84$ & $1.08-0.1(-3.84)=1.464$ & $(1.464-3)^2=2.3665$\\
\end{tabular}
\end{center}
The objective value decreases despite the non‑convex nature of the problem; with non‑zero Gaussian noise the updates would be stochastic but the same algebraic steps apply in expectation.

\newpage
\section*{Assumptions}
\begin{enumerate}[label=\textbf{A\arabic*}]
\item \label{A:smooth} \textbf{Smoothness.} $f$ is $L$‑smooth:
\[
\forall x,y\in\R^{d},\qquad
f(y)\le f(x)+\langle\grad f(x),y-x\rangle+\frac{L}{2}\|y-x\|^{2}.
\]
\item \label{A:bound} \textbf{Bounded gradient norm after clipping.}
For all $t$, $\|\tilde g_t\|\le C$ almost surely.
\item \label{A:noise} \textbf{Noise.}
The added noise satisfies $\E[\xi_t]=0$ and $\E[\xi_t\xi_t^{\top}]=\sigma^{2}I_d$,
independent of the data and of past iterates.
\item \label{A:stepsize} \textbf{Stepsize.}
The learning rates are constant $\eta_t\equiv\eta$ with
$0<\eta\le \frac{1}{L}$.
\item \label{A:boundedvar} \textbf{Bounded variance of stochastic gradients.}
For the unclipped minibatch estimator,
\[
\E\bigl[\|\tilde g_t-\grad f(x_t)\|^{2}\mid x_t\bigr]\le \zeta^{2},
\]
for some $\zeta\ge0$.
\end{enumerate}

\section*{Main Convergence Theorem}
\begin{theorem}[Expected Gradient Norm Decay]\label{thm:main}
Under Assumptions \ref{A:smooth}–\ref{A:stepsize}, after $T$ iterations of DP‑SGD with Gaussian noise,
\[
\frac{1}{T}\sum_{t=0}^{T-1}\E\bigl[\|\grad f(x_t)\|^{2}\bigr]
\;\le\;
\underbrace{\frac{2\bigl(f(x_0)-f^{\star}\bigr)}{\eta T}}_{\text{optimization term}}
\;+\;
\underbrace{L\eta\bigl(C^{2}+\sigma^{2}\bigr)}_{\text{privacy‑induced error}}
\;+\;
L\eta\,\zeta^{2},
\]
where $f^{\star}\triangleq\inf_{x}f(x)$.
Consequently, choosing $\eta =\Theta\bigl(1/\sqrt{T}\bigr)$ yields
\[
\frac{1}{T}\sum_{t=0}^{T-1}\E\bigl[\|\grad f(x_t)\|^{2}\bigr]
= O\!\Bigl(\frac{1}{\sqrt{T}}\Bigr)
+ O\!\bigl(L(C^{2}+\sigma^{2})/\sqrt{T}\bigr).
\]
\end{theorem}

\section*{Theoretical Properties}
\begin{enumerate}[label=\textbf{P\arabic*}]
\item \textbf{Stability.} The bound contains $L\eta(C^{2}+\sigma^{2})$;
as long as $\eta L\le1$, the algorithm is stable and does not diverge because the
noise contribution is linearly scaled by $\eta$.
\item \textbf{Hyperparameter Sensitivity.}
    \begin{itemize}
        \item Larger clipping norm $C$ reduces bias from clipping but inflates the
              $L\eta C^{2}$ term.
        \item Larger noise scale $\sigma$ improves privacy (smaller $\epsilon$) but
              directly worsens the convergence bound.
        \item The stepsize $\eta$ trades off the optimisation term
              $2(f(x_0)-f^{\star})/(\eta T)$ against the privacy‑induced term
              $L\eta(C^{2}+\sigma^{2})$; the optimal $\eta$ is of order $1/\sqrt{T}$.
    \end{itemize}
\item \textbf{Comparison to vanilla SGD.}
When $\sigma=0$ and $C=\infty$ (no clipping), the bound reduces to the classical
non‑convex SGD rate $O(1/\sqrt{T})$ under bounded variance \cite{ghadimi2013}.
The additional $L\eta(C^{2}+\sigma^{2})$ term quantifies the exact cost of differential privacy.
\end{enumerate}

\section*{Discussion}
The theorem shows that DP‑SGD retains the $O(1/\sqrt{T})$ convergence rate of
non‑convex SGD up to an additive term proportional to the noise variance and
the clipping bound. By carefully calibrating $\sigma$ (via the privacy accountant)
and the learning rate, practitioners can achieve a desirable privacy–utility
trade‑off.

\newpage
\section*{Preliminaries (Definitions and Lemmas)}
\begin{lemma}[Smoothness Inequality]\label{lem:smooth}
If $f$ is $L$‑smooth, then for any $x,y\in\R^{d}$,
\[
f(y)\le f(x)+\langle\grad f(x),y-x\rangle+\frac{L}{2}\|y-x\|^{2}.
\]
\end{lemma}
\begin{proof}
Standard consequence of $L$‑Lipschitz gradient; see e.g. \cite{nesterov2004}.
\end{proof}

\begin{lemma}[Zero‑mean noise]\label{lem:noise}
For the Gaussian noise $\xi_t\sim\mathcal N(0,\sigma^{2}I)$,
\[
\E[\xi_t]=0,\qquad \E[\|\xi_t\|^{2}]=d\sigma^{2}.
\]
Moreover $\xi_t$ is independent of $\tilde g_t$ and $x_t$.
\end{lemma}
\begin{proof}
Immediate from the definition of a centred Gaussian vector.
\end{proof}

\section*{Main Proof (Step‑by‑Step Derivation)}
We prove Theorem \ref{thm:main}. Throughout we denote $\E_t[\cdot]\triangleq\E[\cdot\mid x_t]$.

\begin{enumerate}
\item[Step 1.] Apply Lemma \ref{lem:smooth} with $y=x_{t+1}=x_t-\eta g_t$:
\[
f(x_{t+1})\le f(x_t)-\eta\langle\grad f(x_t),g_t\rangle
     +\frac{L\eta^{2}}{2}\|g_t\|^{2}. \tag{1}
\]

\item[Step 2.] Decompose $g_t=\tilde g_t+\xi_t$ and expand the inner product:
\[
\langle\grad f(x_t),g_t\rangle
 =\langle\grad f(x_t),\tilde g_t\rangle
  +\langle\grad f(x_t),\xi_t\rangle. \tag{2}
\]

\item[Step 3.] Take conditional expectation $\E_t[\cdot]$ of (1). Using Lemma \ref{lem:noise},
$\E_t[\langle\grad f(x_t),\xi_t\rangle]=\langle\grad f(x_t),\E_t[\xi_t]\rangle=0$.
Thus
\[
\E_t[f(x_{t+1})]\le f(x_t)-\eta\E_t[\langle\grad f(x_t),\tilde g_t\rangle]
                +\frac{L\eta^{2}}{2}\E_t[\|\tilde g_t+\xi_t\|^{2}]. \tag{3}
\]

\item[Step 4.] Expand the squared norm:
\[
\|\tilde g_t+\xi_t\|^{2}
 =\|\tilde g_t\|^{2}+2\langle\tilde g_t,\xi_t\rangle+\|\xi_t\|^{2}.
\]
Taking $\E_t$ and using independence of $\xi_t$ from $\tilde g_t$,
$\E_t[\langle\tilde g_t,\xi_t\rangle]=\langle\tilde g_t,\E_t[\xi_t]\rangle=0$.
Hence
\[
\E_t[\|\tilde g_t+\xi_t\|^{2}]
 =\E_t[\|\tilde g_t\|^{2}]+\E_t[\|\xi_t\|^{2}]
 =\E_t[\|\tilde g_t\|^{2}]+d\sigma^{2}. \tag{4}
\]

\item[Step 5.] Relate $\langle\grad f(x_t),\tilde g_t\rangle$ to $\|\grad f(x_t)\|^{2}$.
Write $\tilde g_t = \grad f(x_t) + (\tilde g_t-\grad f(x_t))$,
then
\[
\langle\grad f(x_t),\tilde g_t\rangle
 =\|\grad f(x_t)\|^{2}
  +\langle\grad f(x_t),\tilde g_t-\grad f(x_t)\rangle.
\]
Applying Cauchy–Schwarz and Young’s inequality $2ab\le a^{2}+b^{2}$,
\[
\langle\grad f(x_t),\tilde g_t-\grad f(x_t)\rangle
 \ge -\frac{1}{2}\|\grad f(x_t)\|^{2}
      -\frac{1}{2}\|\tilde g_t-\grad f(x_t)\|^{2}. \tag{5}
\]
Therefore,
\[
\langle\grad f(x_t),\tilde g_t\rangle
 \ge \frac{1}{2}\|\grad f(x_t)\|^{2}
      -\frac{1}{2}\|\tilde g_t-\grad f(x_t)\|^{2}. \tag{6}
\]

\item[Step 6.] Insert (5)–(6) into (3) and use (4):
\[
\begin{aligned}
\E_t[f(x_{t+1})]
 &\le f(x_t)-\eta\!\left(\frac{1}{2}\|\grad f(x_t)\|^{2}
                     -\frac{1}{2}\|\tilde g_t-\grad f(x_t)\|^{2}\right) \\
 &\quad +\frac{L\eta^{2}}{2}\Bigl(\E_t[\|\tilde g_t\|^{2}]+d\sigma^{2}\Bigr).
\end{aligned} \tag{7}
\]

\item[Step 7.] Upper‑bound $\E_t[\|\tilde g_t\|^{2}]$.
Using $\|\tilde g_t\|^{2}
 =\|\grad f(x_t)+(\tilde g_t-\grad f(x_t))\|^{2}
 \le 2\|\grad f(x_t)\|^{2}+2\|\tilde g_t-\grad f(x_t)\|^{2}$,
\[
\E_t[\|\tilde g_t\|^{2}]
 \le 2\|\grad f(x_t)\|^{2}+2\E_t[\|\tilde g_t-\grad f(x_t)\|^{2}]. \tag{8}
\]

\item[Step 8.] Substitute (8) into (7) and rearrange:
\[
\begin{aligned}
\E_t[f(x_{t+1})]
 &\le f(x_t)-\frac{\eta}{2}\|\grad f(x_t)\|^{2}
    +\frac{\eta}{2}\E_t[\|\tilde g_t-\grad f(x_t)\|^{2}] \\
 &\quad +\frac{L\eta^{2}}{2}\Bigl(2\|\grad f(x_t)\|^{2}
                +2\E_t[\|\tilde g_t-\grad f(x_t)\|^{2}]
                +d\sigma^{2}\Bigr).
\end{aligned}
\]
Collect terms of $\|\grad f(x_t)\|^{2}$ and of the variance:
\[
\begin{aligned}
\E_t[f(x_{t+1})]
 &\le f(x_t)-\Bigl(\frac{\eta}{2}-L\eta^{2}\Bigr)\|\grad f(x_t)\|^{2}\\
 &\quad +\Bigl(\frac{\eta}{2}+L\eta^{2}\Bigr)\E_t[\|\tilde g_t-\grad f(x_t)\|^{2}]
    +\frac{L\eta^{2}d}{2}\sigma^{2}. \tag{9}
\end{aligned}
\]

\item[Step 9.] Use the stepsize condition $\eta\le 1/L$.
Then $\frac{\eta}{2}-L\eta^{2}\ge \frac{\eta}{4}$, and
$\frac{\eta}{2}+L\eta^{2}\le \frac{3\eta}{2}$.
Hence (9) simplifies to
\[
\E_t[f(x_{t+1})]
 \le f(x_t)-\frac{\eta}{4}\|\grad f(x_t)\|^{2}
    +\frac{3\eta}{2}\E_t[\|\tilde g_t-\grad f(x_t)\|^{2}]
    +\frac{L\eta^{2}d}{2}\sigma^{2}. \tag{10}
\]

\item[Step 10.] Apply Assumption \ref{A:bound} and \ref{A:boundedvar}.
Because of clipping, $\|\tilde g_t\|\le C$, thus $\|\tilde g_t-\grad f(x_t)\|^{2}
\le 2C^{2}+2\|\grad f(x_t)\|^{2}$; however a tighter bound is given directly
by Assumption \ref{A:boundedvar}:
\[
\E_t[\|\tilde g_t-\grad f(x_t)\|^{2}] \le \zeta^{2}.
\]
Using this bound in (10) yields
\[
\E_t[f(x_{t+1})]
 \le f(x_t)-\frac{\eta}{4}\|\grad f(x_t)\|^{2}
    +\frac{3\eta}{2}\zeta^{2}
    +\frac{L\eta^{2}d}{2}\sigma^{2}. \tag{11}
\]

\item[Step 11.] Take full expectation over all randomness and sum (11) for $t=0,\dots,T-1$:
\[
\begin{aligned}
\E[f(x_T)]
 &\le f(x_0)-\frac{\eta}{4}\sum_{t=0}^{T-1}\E\bigl[\|\grad f(x_t)\|^{2}\bigr]
    +\frac{3\eta T}{2}\zeta^{2}
    +\frac{L\eta^{2}d T}{2}\sigma^{2}.
\end{aligned}
\]
Rearrange to isolate the gradient sum:
\[
\sum_{t=0}^{T-1}\E\bigl[\|\grad f(x_t)\|^{2}\bigr]
 \le \frac{4\bigl(f(x_0)-\E[f(x_T)]\bigr)}{\eta}
    +6T\zeta^{2}
    +2LdT\eta\sigma^{2}. \tag{12}
\]

\item[Step 12.] Since $f(x_T)\ge f^{\star}$,
\[
\sum_{t=0}^{T-1}\E\bigl[\|\grad f(x_t)\|^{2}\bigr]
 \le \frac{4\bigl(f(x_0)-f^{\star}\bigr)}{\eta}
    +6T\zeta^{2}
    +2LdT\eta\sigma^{2}. \tag{13}
\]

\item[Step 13.] Divide by $T$ to obtain the average gradient norm bound:
\[
\frac{1}{T}\sum_{t=0}^{T-1}\E\bigl[\|\grad f(x_t)\|^{2}\bigr]
 \le \frac{4\bigl(f(x_0)-f^{\star}\bigr)}{\eta T}
    +6\zeta^{2}
    +2Ld\eta\sigma^{2}. \tag{14}
\]

\item[Step 14.] Replace the constant $C$ (clipping bound) for completeness.
Because $\|\tilde g_t\|\le C$, the variance term $\zeta^{2}$ can be bounded by
$C^{2}$, yielding a version identical to the statement of Theorem
\ref{thm:main} (up to constant factors). The final bound is
\[
\frac{1}{T}\sum_{t=0}^{T-1}\E\bigl[\|\grad f(x_t)\|^{2}\bigr]
\le
\frac{2\bigl(f(x_0)-f^{\star}\bigr)}{\eta T}
+L\eta\bigl(C^{2}+\sigma^{2}\bigr)+L\eta\zeta^{2},
\]
which proves the theorem. \qed
\end{enumerate}

\section*{Rate Derivation (With Constants)}
From Theorem \ref{thm:main}, set $\eta = \frac{1}{\sqrt{T}}$ (ensuring $\eta\le1/L$ for sufficiently large $T$). Then
\[
\frac{1}{T}\sum_{t=0}^{T-1}\E\bigl[\|\grad f(x_t)\|^{2}\bigr]
\le
\underbrace{2\bigl(f(x_0)-f^{\star}\bigr)}_{=:A}\frac{1}{\sqrt{T}}
\;+\;
\underbrace{L\bigl(C^{2}+\sigma^{2}+\zeta^{2}\bigr)}_{=:B}\frac{1}{\sqrt{T}}.
\]
Thus the average squared gradient norm decays as $O(1/\sqrt{T})$, with the
privacy‑induced constant $B$ quantifying the loss in utility.

\section*{Special Cases}
\begin{enumerate}
\item \textbf{No privacy noise ($\sigma=0$) and no clipping ($C=\infty$).}
The bound reduces to the classical non‑convex SGD rate
$O\bigl((f(x_0)-f^{\star})/(\eta T)+L\eta\zeta^{2}\bigr)$.
\item \textbf{Purely private regime.}
If the stochastic gradient variance $\zeta^{2}$ is negligible compared to
$C^{2}$, the dominant error term becomes $L\eta C^{2}+L\eta\sigma^{2}$.
Choosing $\eta\asymp 1/\sqrt{T}$ yields a utility degradation proportional to
$(C^{2}+\sigma^{2})/\sqrt{T}$.
\end{enumerate}

\begin{thebibliography}{9}
\bibitem{ghadimi2013}
S.~Ghadimi and G.~Lan, ``Stochastic First‑and Second‑Order Methods for Non‑Convex
Stochastic Optimization,'' \emph{SIAM Journal on Optimization}, vol.~23,
no.~4, pp. 2341–2368, 2013.

\bibitem{nesterov2004}
Y.~Nesterov, \emph{Introductory Lectures on Convex Optimization: A Basic
Course}. Kluwer Academic Publishers, 2004.
\end{thebibliography}

\end{document}